%%%%%%%%%%%%%%%%%%%%%%%%%%%%%%%%%%%%%%%%%%%%%%%%%%%%%%%%%%%%%%%%%%%%%%%%%%%
%%  chapter1.tex  –  Глава 1: Обзор литературы и анализ пробелов
%%  v9 – Реструктурировано в соответствии с поэтапным охватом:
%%       НВ динамика → распределённость → эффективность → нестационарность.
%%       Покрыты все 12 пересечений 3 условий × 4 требований.
%%       Выявлены 7 пробелов. Связующие утверждения между разделами.
%%       Только заголовки \section добавлены в оглавление.
%%%%%%%%%%%%%%%%%%%%%%%%%%%%%%%%%%%%%%%%%%%%%%%%%%%%%%%%%%%%%%%%%%%%%%%%%%%

\chapter{Анализ динамических моделей графовых нейронных сетей, структурные проблемы и теоретические основы}
\label{ch:literature}

В настоящей главе представлен обзор существующей литературы по динамическим графовым нейронным сетям и проблемам их робастности, следующий поэтапному охвату диссертации: от динамики непрерывного времени на эволюционирующих графах, через распределённое федеративное обучение с приватностью, к вычислительной эффективности для крупномасштабных графов и, наконец, к нестационарной адаптации с оценкой неопределённости. Каждый раздел адресует одно или несколько пересечений трёх рабочих условий (состязательная среда, распределённая динамика, нестационарность) и четырёх сквозных требований (робастность, точность, эффективность, приватность), выстраивая систематическое перечисление семи исследовательских пробелов, мотивирующих методы, разработанные в последующей главе.


%%=====================================================================
\section{Графовые нейронные сети: архитектура, уязвимость и проблема двойного входа}
\label{sec:gnn_vulnerability}
%%=====================================================================

\subsection{Архитектура передачи сообщений и агрегация окрестностей}

Графовые нейронные сети~\cite{Kipf2017,Velickovic2018,Hamilton2017} расширяют глубокое обучение на неевклидовы реляционные данные путём итеративной агрегации информации от соседей узла. Каноническая формулировка передачи сообщений обновляет представление узла $v$ на слое $k+1$ как
\begin{equation}
\bh_v^{(k+1)}
= \phi\!\Bigl(\bh_v^{(k)},\;
  \bigoplus_{u\in\calN(v)}
    \psi\!\bigl(\bh_u^{(k)},\, \mathbf{e}_{uv}\bigr)\Bigr),
\label{eq:ch1_mp}
\end{equation}
где $\phi$ и $\psi$ --- обучаемые функции, $\bigoplus$ --- перестановочно-инвариантный агрегатор (сумма, среднее или максимум), а $\calN(v)$ обозначает окрестность $v$. Данная схема охватывает GCN~\cite{Kipf2017}, GAT~\cite{Velickovic2018}, GraphSAGE~\cite{Hamilton2017} и их темпоральные расширения~\cite{Xu2020tgat,Rossi2020,Pei2020tgn,Kazemi2020,Schlichtkrull2018}. Выразительность парадигмы передачи сообщений формально охарактеризована через иерархию изоморфизмов Вейсфейлера--Лемана~\cite{Xu2019gin}, устанавливающую фундаментальные ограничения на то, какие свойства уровня графа эти архитектуры могут различать.

Графовые нейронные сети применяются в здравоохранении (предсказание свойств молекул на графах атомов и связей, графы взаимодействия пациентов для рекомендаций по лечению), финансах (анализ графов транзакций для обнаружения мошенничества и противодействия отмыванию денег), кибербезопасности (графы сетевого трафика для обнаружения вторжений), промышленных системах (графы слияния датчиков для мониторинга SCADA и АСУ ТП) и автономных системах (графы слияния датчиков для восприятия)~\cite{Goodfellow2014gan,Gilmer2017,Gilmer2018,Ying2021,Wu2021,Luo2021}. Широта этих применений подчёркивает актуальность анализа робастности.


\subsection{Уязвимость двойного входа}

В отличие от обычных нейронных сетей, принимающих единый тензор данных, графовая нейронная сеть одновременно получает два входа: числовые \emph{входные данные} $\mathbf{X}\in\mathbb{R}^{n\times d}$ (матрица признаков, кодирующая атрибуты сущностей) и реляционную \emph{входную структуру} $\mathbf{A}\in\{0,1\}^{n\times n}$ (матрица смежности, кодирующая попарные связи). Сеть формирует выход путём распространения данных вдоль структуры через итеративную агрегацию окрестностей, и возмущение \emph{любого} из входов --- данных или структуры --- может нарушить выход~\cite{Zuegner2018,Bojchevski2019,Dai2018}.

Z\"{u}gner и др.~\cite{Zuegner2018} продемонстрировали, что изменение всего пяти рёбер в графе цитирования может переключить предсказанную метку целевого узла, в то время как Bojchevski и G\"{u}nnemann~\cite{Bojchevski2019} показали, что глобальные атаки отравления снижают точность классификации узлов до 20\% на бенчмарк-наборах данных. Dai и др.~\cite{Dai2018} расширили эти результаты на возмущение графов на основе обучения с подкреплением, а последующие работы~\cite{Xu2019b,Wu2019} подтвердили, что как признаковые, так и топологические возмущения переносятся между архитектурами ГНС. Biggio и Roli~\cite{Biggio2018} предоставили исчерпывающую таксономию состязательных атак на машинное обучение, а Carlini и Wagner~\cite{Carlini2017} установили эталонные методологии оценки. Kurakin и др.~\cite{Kurakin2017} продемонстрировали, что состязательные возмущения сохраняются в условиях реального мира, а Papernot и др.~\cite{Papernot2016,Papernot2017pate} показали осуществимость атак чёрного ящика через переносимость. Athalye и др.~\cite{Athalye2018} впоследствии показали, что обфусцированные градиенты создают ложное чувство безопасности, мотивируя необходимость сертифицированных, а не эмпирических защит. Bhagoji и др.~\cite{Bhagoji2019} проанализировали состязательные атаки в распределённых условиях, Jia и Liang~\cite{Jia2019} исследовали отравление данных для рекомендательных систем, а Bose и Hamilton~\cite{Bose2020} разработали состязательные атаки на графовые вложения. Sun и др.~\cite{Sun2022gnn_robustness} обозрели методы робастности, специфичные для графов, а Jin и др.~\cite{Jin2020} разработали системы состязательного обучения графов. Компромисс между состязательной робастностью и точностью был формально охарактеризован Tsipras и др.~\cite{Tsipras2019}, а Tram\`{e}r и др.~\cite{Tramer2018} продемонстрировали ансамблевые атаки, преодолевающие маскировку градиентов.

Gama, Bruna и Ribeiro~\cite{Gama2020} предложили систему устойчивости, ограничивающую чувствительность выхода графовых фильтров через константу Липшица их частотной характеристики. Однако эти границы быстро растут с глубиной сети и не распространяются на темпоральные, распределительные или состязательные условия, с которыми сталкиваются реальные графовые системы.


\subsection{Задокументированные отказы в критически важных областях}

Уязвимость графовых нейронных сетей не является теоретической проблемой --- она задокументирована в рецензируемых исследованиях в четырёх критически важных областях. В здравоохранении градиентные атаки на чувствительность рёбер на графовые изоморфные сети вызвали падение AUC приблизительно на 25\% на бенчмарках MoleculeNet (Tox21, ClinTox, SIDER), при этом токсичные соединения были классифицированы как безопасные~\cite{EdgeSensitivity2025}. В финансах атака MonTi с внедрением группы почти утроила ошибку классификации ГНС-детекторов мошенничества с 25,7\% до 68,6\%~\cite{MonTi2025}. В промышленном IoT отравление данных федеративных детекторов вторжений вызвало коллапс точности на 64 процентных пункта с 94,2\% до 30,0\% при не-IID условиях~\cite{FerragPoisoning2023}. В кибербезопасности атака PROVNINJA снизила обнаружение ГНС-систем на основе провенанса до 59\%, сделав APT-атакующих невидимыми~\cite{PROVNINJA2023}. Эти задокументированные отказы подтверждают, что робастность ГНС отстаёт от точности предсказаний на 25--65 процентных пунктов во всех критически важных областях.

\medskip
\noindent\textit{Итоги раздела.} Механизм передачи сообщений, придающий графовым нейронным сетям выразительную способность, также создаёт поверхность атаки двойного входа, эксплуатируемую как через признаковые, так и через топологические возмущения. Существующие границы устойчивости не распространяются на темпоральные, распределительные или состязательные условия. В следующем разделе рассматривается, могут ли формулировки непрерывного времени обеспечить недостающие сертифицированные гарантии робастности.


%%=====================================================================
\section{Динамика непрерывного времени на эволюционирующих графах}
\label{sec:ct_dynamics}
%%=====================================================================

\subsection{Нейронные обыкновенные дифференциальные уравнения для представления графов}

Chen и др.~\cite{Chen2018} предложили нейронные обыкновенные дифференциальные уравнения (нейронные ОДУ), заменяющие дискретные послойные обновления обычных сетей параметризацией непрерывной глубины:
\begin{equation}
\frac{d\bh(t)}{dt} = f_\theta\!\bigl(\bh(t),\, t\bigr),
\label{eq:ch1_ode}
\end{equation}
где $f_\theta$ --- нейронная сеть, параметризующая векторное поле. Выход в любой момент времени $T$ получается численным интегрированием уравнения~\eqref{eq:ch1_ode} от $t=0$ до $t=T$ с помощью решателей с адаптивным шагом, а градиенты вычисляются с эффективным использованием памяти через метод сопряжённых чувствительностей~\cite{Pontryagin1962}. Данная формулировка была расширена на графовые данные: Chamberlain и др.~\cite{Chamberlain2021} сформулировали графовую нейронную диффузию как непрерывное УЧП на графе, а Rusch и др.~\cite{Rusch2022} разработали связанные осцилляторные графовые нейронные ОДУ, предотвращающие чрезмерное сглаживание. Rubanova и др.~\cite{Rubanova2019} расширили нейронные ОДУ на нерегулярно дискретизированные временные ряды, Dupont и др.~\cite{Dupont2019} предложили расширенные нейронные ОДУ с пространствами состояний повышенной размерности, а Kidger и др.~\cite{Kidger2020} разработали нейронные управляемые дифференциальные уравнения для нерегулярных временных рядов. Grathwohl и др.~\cite{Grathwohl2019} связали нейронные ОДУ с непрерывными нормализующими потоками, а Finlay и др.~\cite{Finlay2020} предложили методы регуляризации для стабильного обучения. Hasani и др.~\cite{Hasani2021} разработали нейронные ОДУ с жидкими постоянными времени для адаптивных вычислений, а Li и др.~\cite{Li2020ode} применили нейронные ОДУ к непрерывным графовым нейронным сетям на статических топологиях.

Темпоральные графовые сети, такие как TGAT~\cite{Xu2020tgat} и TGN~\cite{Rossi2020}, моделируют взаимодействия с временными метками на эволюционирующих графах, но работают в дискретном времени --- каждое событие вызывает один шаг обновления. Это принудительно привязывает нерегулярные события графа к фиксированной вычислительной тактовой сетке, теряя информацию о точном времени и порядке взаимодействий между шагами.


\subsection{Устойчивость по Липшицу и граница Грёнуолла}

Ключевое преимущество формулировок непрерывного времени --- их присущая гладкость. Поскольку представления эволюционируют согласно дифференциальному уравнению с ограниченной константой Липшица $L_g$, неравенство Грёнуолла обеспечивает замкнутый сертификат робастности:
\begin{equation}
\|\mathbf{h}_1(T) - \mathbf{h}_2(T)\| \leq e^{L_g T}\|\mathbf{x}_1 - \mathbf{x}_2\|,
\label{eq:ch1_gronwall}
\end{equation}
где $\mathbf{h}_i(T)$ --- скрытое состояние траектории $i$ в момент $T$, а $\mathbf{x}_i$ --- начальный вход. Эта граница утверждает, что малые возмущения входа могут вызвать лишь ограниченные, гладкие изменения выхода --- гладкость управляющего потока математически ограничивает степень изменения предсказания злоумышленником. Графовые нейронные сети дискретного времени не имеют такого непрерывного механизма, поэтому возмущения непредсказуемо накапливаются от слоя к слою~\cite{Zhao2023hamilton}.

Рандомизированное сглаживание~\cite{Cohen2019} и распространение интервальных границ~\cite{Gowal2019} обеспечивают альтернативные пути сертификации, но они были разработаны для классификаторов изображений и не учитывают комбинаторную структуру матриц смежности графов или интегрирование нейронных ОДУ в непрерывном времени. Lecuyer и др.~\cite{Lecuyer2019} связали дифференциальную приватность с сертифицированной робастностью, а Salman и др.~\cite{Salman2019} объединили рандомизированное сглаживание с состязательным обучением. Shafahi и др.~\cite{Shafahi2019} разработали свободное состязательное обучение для повышения эффективности, а Wong и Kolter~\cite{Wong2020} предоставили масштабируемое верифицированное обучение через выпуклую релаксацию.


\subsection{Ограничения существующих подходов непрерывного времени}

Несмотря на теоретическую перспективность, ни одна существующая работа не выводит сертифицированные границы робастности из гладкости управляющего дифференциального уравнения на темпорально эволюционирующих топологиях графов. Существующие модели графов непрерывного времени либо работают на статических графах~\cite{Chamberlain2021}, либо обрабатывают темпоральные рёбра без анализа Липшица~\cite{Xu2020tgat,Rossi2020,Pei2020tgn,Kazemi2020,Wang2021tedic}. Многомасштабная темпоральная динамика (захватывающая как субсекундные переходные процессы, так и часовые структурные тенденции) не была решена в рамках сертифицированной системы. Более того, обучение с $O(1)$-памятью через сопряжённый метод не было сочетано с темпоральной адаптивной нормализацией для эволюционирующих топологий графов.

\medskip
\noindent\textbf{Пробел~1:} \emph{Отсутствие сертифицированной робастности при динамике графов непрерывного времени.} Ни одна существующая архитектура не обеспечивает формально сертифицированные границы робастности для графовых нейронных сетей, работающих в непрерывном времени на темпорально эволюционирующих топологиях, сочетая анализ Липшица управляющего векторного поля с обучением через сопряжённый метод с эффективным использованием памяти и многомасштабной темпоральной динамикой.

\medskip
\noindent\textit{Итоги раздела.} Формулировки непрерывного времени предлагают принципиальный путь к сертифицированной робастности через динамику с ограниченной константой Липшица, но существующие подходы не распространяются на эволюционирующие графы с темпоральными потоками событий. В следующем разделе рассматривается, как эти модели должны быть адаптированы для распределённых сред, в которых данные не могут быть централизованы.


%%=====================================================================
\section{Распределённое обучение на графах с гарантиями приватности}
\label{sec:distributed}
%%=====================================================================

\subsection{Основы федеративного обучения}

McMahan и др.~\cite{McMahan2017} предложили федеративное усреднение (FedAvg), позволяющее проводить распределённое обучение на множестве участников без централизации сырых данных. Каждый участник вычисляет локальные обновления градиентов на своём фрагменте данных, а центральный сервер агрегирует эти обновления в глобальную модель. Данная парадигма необходима для графовых данных в здравоохранении (графы взаимодействия пациентов между больницами), финансах (графы транзакций между банками) и кибербезопасности (графы трафика между центрами обработки данных), где регуляции приватности запрещают обмен данными~\cite{Kairouz2021,Konecny2016,Mothukuri2021}. Dwork~\cite{Dwork2014} заложил основы алгоритмической приватности, а McMahan и др.~\cite{McMahan2018dp} объединили федеративное обучение с дифференциальной приватностью в масштабе. Li и др.~\cite{Li2020fedprox} решили проблему гетерогенного федеративного обучения через проксимальную регуляризацию, а Zhao и др.~\cite{Zhao2020fl_noniid} охарактеризовали проблему не-IID в федеративных условиях. Wang и др.~\cite{Wang2020fedma} предложили федеративное согласованное усреднение для гетерогенных архитектур. Sun и Lyu~\cite{Sun2019fl_poison} продемонстрировали атаки отравления на федеративное обучение, Fang и др.~\cite{Fang2020} разработали атаки отравления локальных моделей, а El-Mhamdi и др.~\cite{ElMhamdi2018} установили информационно-теоретические пределы византийско-устойчивой агрегации. Geyer и др.~\cite{Geyer2017} предложили дифференциальную приватность на уровне клиентов для федеративного обучения, а Agarwal и др.~\cite{Agarwal2018cpsgd,Bernstein2017,Bernstein2019} разработали коммуникационно-эффективный распределённый SGD с гарантиями приватности. Wei и др.~\cite{Wei2020dpfl} проанализировали сходимость дифференциально приватного федеративного обучения.


\subsection{Византийско-устойчивая агрегация}

Стандартное федеративное усреднение уязвимо к византийским участникам, отправляющим повреждённые обновления градиентов. Blanchard и др.~\cite{Blanchard2017} предложили Krum, который выбирает обновление, наиболее близкое к его соседям, в то время как Yin и др.~\cite{Yin2018} показали, что покоординатные усечённые средние достигают минимаксно-оптимальных скоростей при византийском повреждении. Для графовых данных сама топология может быть частично контролируема противником, добавляя комбинаторную поверхность атаки, отсутствующую в стандартных федеративных условиях.


\subsection{Дифференциальная приватность и оптимальный транспорт}

Дифференциальная приватность~\cite{Dwork2006} формализует конфиденциальность на уровне записей данных через условие
\begin{equation}
\Prob\bigl[\calM(\calD)\in S\bigr]
\;\le\;
e^{\epsilon}\;\Prob\bigl[\calM(\calD')\in S\bigr] + \delta,
\end{equation}
а композиция через моменты-бухгалтер~\cite{Abadi2016} позволяет точно отслеживать кумулятивный расход приватности по раундам обучения. Оптимальный транспорт~\cite{Villani2009,Cuturi2013} обеспечивает геометрию на распределениях вероятностей через расстояние Вассерштейна, предлагая принципиальный механизм для выравнивания гетерогенных распределений фрагментов графов между участниками. Courty и др.~\cite{Courty2017} применили оптимальный транспорт к адаптации доменов, Flamary и др.~\cite{Flamary2021pot} выпустили библиотеку Python Optimal Transport для масштабируемых вычислений, а Genevay и др.~\cite{Genevay2018} разработали выборочно-эффективные дивергенции Синкхорна для генеративного моделирования. Однако ни один существующий алгоритм не гарантирует совместно византийскую устойчивость, дифференциальную приватность и выравнивание распределений по оптимальному транспорту для графовых данных.

\medskip
\noindent\textbf{Пробел~2:} \emph{Неблагоприятный компромисс приватность--качество в федеративных графах.} Ни один существующий алгоритм федеративного обучения не обеспечивает совместно византийско-устойчивую агрегацию, формальную $(\epsilon,\delta)$-дифференциальную приватность и выравнивание распределений по оптимальному транспорту Вассерштейна для графовых данных, где сама топология может быть частично контролируема противником.

\medskip
\noindent\textit{Итоги раздела.} Распределённое обучение на графах требует одновременной устойчивости к повреждениям, защиты приватности и выравнивания распределений --- комбинации, которой ни один существующий метод не достигает. В следующем разделе рассматриваются проблемы вычислительной эффективности, которые должно преодолеть любое развёрнутое решение.


%%=====================================================================
\section{Вычислительная эффективность для крупномасштабных темпоральных графов}
\label{sec:efficiency}
%%=====================================================================

\subsection{Квадратичное узкое место графовых трансформеров}

Архитектуры графовых трансформеров~\cite{Vaswani2017,Ying2021} достигают высоких результатов на задачах уровня графа, позволяя каждому узлу обращать внимание на каждый другой узел. Однако этот механизм самовнимания масштабируется квадратично ($\mathcal{O}(L^2)$) с числом узлов или событий, делая его непомерно дорогим для крупномасштабных темпоральных графов с миллионами взаимодействий. Варианты с линейным вниманием~\cite{Katharopoulos2020} снижают это до $\mathcal{O}(L)$, но жертвуют выразительностью, необходимой для графо-топологического рассуждения.


\subsection{Структурированные модели пространства состояний}

Gu и др.~\cite{Gu2022} предложили структурированные модели пространства состояний (S4), переформулирующие последовательную обработку как линейную динамическую систему, свёрточное ядро которой вычисляется за время $\mathcal{O}(L\log L)$ через быстрое преобразование Фурье. Селективный вариант (Mamba)~\cite{Dao2024} добавляет входозависимое стробирование для улучшенного моделирования дискретных токенов. Behrouz и др.~\cite{Behrouz2024graphmamba} начали применять эти модели к графовым данным (GraphMamba), но их робастность при целенаправленном возмущении и способность формировать надёжные оценки уверенности не исследовались. Теоретические основы моделей пространства состояний восходят к системе HiPPO~\cite{Gu2020hippo} для оптимальной полиномиальной проекции, а Orvieto и др.~\cite{Orvieto2023} предоставили унифицированный анализ линейных рекуррентных архитектур. Poli и др.~\cite{Poli2023} разработали операторы Hyena как свёрточные альтернативы, а Choromanski и др.~\cite{Choromanski2021,Perez2021} предложили внимание на случайных признаках для эффективных трансформеров. Dosovitskiy и др.~\cite{Dosovitskiy2021,Salvi2024} продемонстрировали, что трансформеры зрения достигают конкурентной производительности через токенизацию на основе патчей, установив парадигму, впоследствии адаптированную к графовым последовательностям.


\subsection{Многоуровневые представления}

Темпоральные графы, в которых узлы и рёбра несут разные семантические типы и эволюционируют на разных временных масштабах, требуют представлений, захватывающих структуру на нескольких уровнях детализации. Существующие темпоральные ГНС обучают единое, унифицированное вложение на узел, которое не способно одновременно разрешить паттерны уровня сервисов, зависимости уровня трассировок и динамику уровня узлов. Многоуровневые представления дороги в вычислении и не были проанализированы на устойчивость при топологической эволюции или на их взаимодействие с катастрофическим забыванием.

\medskip
\noindent\textbf{Пробел~3:} \emph{Ограниченная кросс-распределительная обобщаемость на темпоральных графах.} Ни одна существующая архитектура не обучает многоуровневые темпоральные представления графов (сервис, трассировка, узел) с одновременным предотвращением катастрофического забывания при эволюции топологии и достижением вычислительной экономии через структурированную разреженность.

\medskip
\noindent\textbf{Пробел~4:} \emph{Ограничения дискретного времени и квадратичная стоимость.} Структурированные модели пространства состояний, предлагающие близкую к линейной стоимость, не были адаптированы к графовым данным с формальными гарантиями робастности, прогрессивной состязательной дистилляцией или PAC-байесовскими сертификатами обобщения.

\medskip
\noindent\textit{Итоги раздела.} Эффективность и многоуровневое представление ставят взаимосвязанные задачи: более богатые представления требуют больше вычислений, а эффективные архитектуры лишены гарантий робастности. В следующем разделе рассматривается, как развёрнутые модели должны адаптироваться к нестационарным условиям.


%%=====================================================================
\section{Нестационарная адаптация и оценка неопределённости}
\label{sec:nonstationarity}
%%=====================================================================

\subsection{Дрейф распределения в развёрнутых моделях}

Реальные графовые данные дрейфуют после развёртывания: топологии эволюционируют, паттерны атак меняются, появляются новые протоколы, и статические модели незаметно теряют точность без предупреждения. Это недопустимо в критически важных системах, где устаревшая модель является опасной моделью. Дрейф концепта был обширно исследован в табличных и потоковых условиях~\cite{Gama2014,Lu2019}, но его взаимодействие с графовыми данными --- где одновременно смещаются как распределения признаков, так и топологические паттерны --- создаёт уникальные проблемы. Parisi и др.~\cite{Parisi2019} обозрели непрерывное обучение на протяжении жизни с нейронными сетями, а Zenke и др.~\cite{Zenke2017,Wilson2020} предложили синаптический интеллект как альтернативу EWC для предотвращения катастрофического забывания.


\subsection{Байесовские подходы к неопределённости}

Байесовские нейронные сети формируют апостериорные предсказательные распределения, декомпозирующие полную неопределённость на эпистемическую (модельную) и алеаторную (данных) компоненты~\cite{Gal2016,Blundell2015,Lakshminarayanan2017,Kendall2017}. Вариационный вывод аппроксимирует неразрешимое апостериорное распределение через параметризованные распределения, обеспечивая оценку неопределённости за счёт дополнительных прямых проходов (дропаут Монте-Карло) или ансамблевого обучения. Однако интеграция вариационного байесовского внимания с оценкой неопределённости на иерархических гауссовских процессах для графовой онлайн-адаптации не была достигнута. Ovadia и др.~\cite{Ovadia2019} провели бенчмарк оценки неопределённости при сдвиге набора данных, обнаружив, что большинство методов значительно деградируют. Dusenberry и др.~\cite{Dusenberry2020} предложили эффективный и масштабируемый байесовский вывод для нейронных сетей, а Foong и др.~\cite{Foong2019} проанализировали выразительность приближённого байесовского вывода в нейронных сетях. Depeweg и др.~\cite{Depeweg2018} формализовали декомпозицию неопределённости на эпистемическую и алеаторную компоненты для принятия решений при модельной неопределённости. Kuleshov и др.~\cite{Kuleshov2018} установили метрики калибровки для современных вероятностных прогнозов.


\subsection{Онлайн-обучение с сублинейным регретом}

Онлайн-выпуклая оптимизация обеспечивает основу для последовательного принятия решений при состязательных потоках данных с алгоритмами типа следование-за-регуляризованным-лидером, достигающими сублинейных границ регрета~\cite{Shalev2012,Arjovsky2017}. Обновление мультипликативными весами достигает регрета $R(T)\leq\sqrt{T\log|S|}$ относительно лучшей фиксированной стратегии ретроспективно. Однако применение этих границ к адаптации графовых нейронных сетей при одновременном распределительном сдвиге и состязательном отравлении остаётся неисследованным.

\medskip
\noindent\textbf{Пробел~5:} \emph{Прямая несовместимость при смене протокольного режима.} Ни одна существующая предметно-ориентированная фундаментальная модель не сочетает обработку на основе пространства состояний, дополненного ОДУ, с многозадачным моделированием графовых последовательностей событий для прямо совместимого обнаружения при эволюционирующих протоколах. Постквантовый криптографический ландшафт, от алгоритма Шора~\cite{Shor1997} через решёточные конструкции~\cite{Ducas2018,Proos2003} и кодовые схемы~\cite{McEliece1978}, вносит фундаментальные изменения в статистические признаки, наблюдаемые на рёбрах графа. Парадигма фундаментальных моделей, установленная Brown и др.~\cite{Brown2020} и Devlin и др.~\cite{Devlin2019}, была расширена на предметно-ориентированные приложения, при этом Touvron и др.~\cite{Touvron2023} продемонстрировали альтернативы с открытыми весами, а Wei и др.~\cite{Wei2022cot} установили возможности рассуждения по цепочке мыслей, обеспечивающие структурированный анализ угроз.

\medskip
\noindent\textbf{Пробел~6:} \emph{Неинтегрированная неопределённость и состязательная робастность.} Ни одна существующая архитектура не сочетает вариационное байесовское внимание с оценкой неопределённости на иерархических гауссовских процессах для графовых данных, обеспечивая одновременно декомпозицию эпистемической и алеаторной неопределённости, онлайновое обнаружение дрейфа и гарантии сублинейного регрета при состязательной нестационарности.

\medskip
\noindent\textit{Итоги раздела.} Нестационарная адаптация требует онлайн-обучения с формальными гарантиями регрета и калиброванной декомпозиции неопределённости --- возможностей, недоступных в существующих графовых архитектурах. Последний пробел касается того, как эти разнообразные гарантии могут быть объединены в единую систему сертификации.


%%=====================================================================
\section{Единая сертификация робастности}
\label{sec:certification}
%%=====================================================================

\subsection{Теоретико-игровые модели угроз}

Tambe~\cite{Tambe2011} и Sinha и др.~\cite{Sinha2018} формализовали распределение ресурсов безопасности как игры Штакельберга, где защитник принимает стратегию, предвосхищая лучший ответ противника. В контексте графовых нейронных сетей защитником является классификатор, а противник возмущает граф для максимизации ошибки классификации. Существующие теоретико-игровые модели для ГНС обычно предполагают однократные взаимодействия с упрощёнными бюджетами возмущений и не учитывают последовательную, адаптивную природу реальных противников или многометодную композиционную структуру современных систем защиты графов. Br\"{u}ckner и Scheffer~\cite{Bruckner2011} формализовали предсказательную игру Штакельберга для состязательной классификации, а Shoham и Leyton-Brown~\cite{Shoham2009} заложили основы мультиагентных систем. Balcan и др.~\cite{Balcan2015} разработали алгоритмы онлайн-обучения со стратегическими взаимодействиями.


\subsection{Существующие методы сертификации}

Рандомизированное сглаживание~\cite{Cohen2019,Salman2019} обеспечивает вероятностные сертификаты робастности для непрерывных входных пространств. Распространение интервальных границ~\cite{Gowal2019,Wong2018} обеспечивает детерминированные границы путём распространения интервально-значных входов через сеть. Однако ни один из этих подходов не был адаптирован к дискретной, комбинаторной структуре матриц смежности графов, и ни одна единая система не обеспечивает доказуемые гарантии одновременно против возмущений признаков, сдвигов временного порядка, модификаций топологии графа и распределительного дрейфа в рамках единой методологии оценки. PAC-байесовская теория~\cite{McAllester1999,Catoni2007} обеспечивает дополнительный путь сертификации через границы обобщения, при этом Dziugaite и Roy~\cite{Dziugaite2017} продемонстрировали невакуумные границы для глубоких сетей, Guedj~\cite{Guedj2019} обозрел современный PAC-байесовский ландшафт, а Letarte и др.~\cite{Letarte2019} связали PAC-байесовские границы с классификаторами на основе голосования большинством.

\medskip
\noindent\textbf{Пробел~7:} \emph{Упрощённые теоретико-игровые модели угроз.} Ни одна существующая единая система сертификации не сочетает теоретико-игровой выбор стратегий (равновесия Штакельберга с онлайн-обучением без регрета), распространение интервальных границ, оценку Липшица на графовых операторах и рандомизированное сглаживание, адаптированное к дискретным графовым структурам, в рамках единой методологии, связывающей все компонентные методы через регуляризацию согласованности.

\medskip
\noindent\textit{Итоги раздела.} Существующие методы сертификации адресуют отдельные типы возмущений изолированно и не распространяются на комбинаторную структуру графов или композиционную архитектуру многометодных систем защиты. Необходима единая система для связывания всех гарантий, разработанных в предыдущих разделах.


%%=====================================================================
\section{Сводка выявленных исследовательских пробелов}
\label{sec:gap_summary}
%%=====================================================================

Систематический анализ, проведённый в разделах~\ref{sec:gnn_vulnerability}--\ref{sec:certification}, выявляет семь исследовательских пробелов, которые в совокупности определяют пространство открытых задач, решаемых в настоящей диссертации. Таблица~\ref{tab:gaps} сопоставляет каждый пробел рабочим условиям и требованиям, которые он затрагивает, и указывает метод, разработанный в Главе~2 для его решения.

\begin{table}[htbp]
\centering
\caption{Семь выявленных исследовательских пробелов, их покрытие условий и требований, а также методы, предложенные для их решения в настоящей диссертации.}
\label{tab:gaps}
\small
\begin{tabular}{p{0.5cm} p{5.5cm} p{2.5cm} p{2.0cm} p{2.5cm}}
\toprule
\textbf{\#} & \textbf{Описание пробела} & \textbf{Условия} & \textbf{Требо\-вания} & \textbf{Предлож. метод} \\
\midrule
1 & Отсутствие сертифиц. робастности при НВ динамике графов & У1, У2 & Т1, Т2 & М1: CT-TGNN \\
\addlinespace
2 & Неблагоприятный компромисс приватность--качество в федерат. графах & У3 & Т1, Т4 & М2: FedLLM-API \\
\addlinespace
3 & Огранич. кросс-распредел. обобщаемость на темпор. графах & У2 & Т2, Т3 & М3: TripleE-TGNN \\
\addlinespace
4 & Ограничения дискретного времени и квадратичная стоимость & У1, У2 & Т1, Т3 & М4: MambaShield \\
\addlinespace
5 & Прямая несовместимость при смене протокольного режима & У2 & Т2, Т3 & CyberSecLLM \\
\addlinespace
6 & Неинтегрированная неопредел. и состязат. робастность & У1, У2 & Т1, Т2 & М5/М6: Стох.\ трансф.\ + UC-HGP \\
\addlinespace
7 & Упрощённые теоретико-игровые модели угроз & У1, У2, У3 & Т1, Т2, Т3, Т4 & М7: Сертификация \\
\bottomrule
\end{tabular}
\end{table}

В совокупности семь пробелов охватывают все три рабочих условия (У1: состязательная среда, У2: нестационарность, У3: распределённая динамика) и все четыре требования (Т1: робастность, Т2: точность, Т3: эффективность, Т4: приватность), подтверждая, что двенадцать пересечений матрицы условие--требование $3\times4$ покрыты предложенными методами. Каждый пробел продемонстрирован как нерешённый существующими работами через анализ, представленный в данной главе.

\medskip
\noindent\textit{Итоги главы.} В настоящей главе заложены теоретические основы и проведён обзор существующей литературы по робастности графовых нейронных сетей, динамике непрерывного времени, распределённому обучению, вычислительной эффективности, нестационарной адаптации и единой сертификации. Выявлены семь конкретных исследовательских пробелов, каждый из которых подкреплён свидетельствами из литературы об отсутствии существующих решений. В следующей главе представлены математические методы, алгоритмы и архитектурные решения, разработанные для заполнения этих пробелов, в том же последовательном порядке: сертифицированная динамика непрерывного времени (Пробел~1), распределённое федеративное обучение (Пробел~2), многоуровневые вложения (Пробел~3), эффективный вывод на основе пространства состояний (Пробел~4), прямо совместимые фундаментальные модели (Пробел~5), адаптация с калиброванной неопределённостью (Пробел~6) и единая теоретико-игровая сертификация (Пробел~7).

